% ======================================================================
% Supplementary Information for
% "Antibody-Antigen Affinity and Mutational ddG Prediction with
%  Chain-Aware Protein Language Modeling"
% Append after \appendix in the main paper, or compile standalone.
% ======================================================================
\documentclass{article}
\usepackage[letterpaper,margin=1in]{geometry}
\usepackage{times}
\usepackage[utf8]{inputenc}\usepackage[T1]{fontenc}
\usepackage{amsmath,amssymb,booktabs,graphicx,xcolor,multirow}
\usepackage[hidelinks]{hyperref}\usepackage{caption}
\captionsetup{font=small,labelfont=bf}
\newcommand{\ddg}{\ensuremath{\Delta\Delta G}}
\newcommand{\pkd}{\ensuremath{pK_d}}
\renewcommand{\thesection}{S\arabic{section}}
\renewcommand{\thetable}{S\arabic{table}}
\renewcommand{\thefigure}{S\arabic{figure}}
\title{\vspace{-2em}\bfseries Supplementary Information:\\ Chain-Aware Antibody--Antigen Affinity and Mutational \ddg{}}
\author{}\date{}
\begin{document}\maketitle\vspace{-2em}

\section{Model and training details}
The encoder is a frozen ESM-2 650M (\texttt{esm2\_t33\_650M\_UR50D}); no ESM-2 parameters are updated.
Heavy, light, and antigen streams are projected to $d{=}256$ (Linear--LayerNorm--GELU, dropout $0.1$).
Heavy--light self-attention uses $8$ heads; the gated cross-attention is stacked twice. For absolute
affinity the model is trained with MSE on the cosine mapped to \pkd{} via the training-fold extrema
$\pkd={\pkd}_{\min}+({\pkd}_{\max}-{\pkd}_{\min})(\hat s+1)/2$. Optimization: AdamW, lr $10^{-4}$, weight decay
$10^{-2}$, batch size $32$, gradient clipping $1.0$, $\leq 50$ epochs, validation-Pearson early stopping
(patience $10$). For \ddg{} the two AbAffinity towers are frozen and weight-shared; only the anchor scale
$s$ (initialised to $RT\ln 10=1.364$), the interface interaction, and the learned interface head are
trained, with a weighted-Huber(hinge) loss ($\delta{=}1,\tau{=}2,\alpha{=}1$) plus a pairwise ranking term
(coefficient $0.5$) and reverse-mutation augmentation. \ddg{} numbers are mean $\pm$ s.d.\ over three seeds
with per-fold $z$-scoring; a full 10-fold \ddg{} run completes in minutes on one GPU.

\section{Absolute-affinity supplementary tables}

\begin{table}[h]\centering
\caption{Architecture comparison on SAAINT-DB (10-fold, mean $\pm$ s.d.). AbAffinity's tri-stream
factorization is best on both splits, with the largest margin on the antigen-cold split.}
\begin{tabular}{llcc}
\toprule
Split & Model & Pearson $r$ & Spearman $\rho$\\
\midrule
\multirow{3}{*}{random}
 & AbAffinity (tri-stream) & $\mathbf{0.842\pm0.022}$ & $\mathbf{0.827\pm0.027}$\\
 & fused two-stream        & $0.815\pm0.028$ & $0.804\pm0.030$\\
 & concat+MLP              & $0.817\pm0.033$ & $0.787\pm0.049$\\
\midrule
\multirow{3}{*}{antigen-cold}
 & AbAffinity (tri-stream) & $\mathbf{0.600\pm0.072}$ & $\mathbf{0.574\pm0.069}$\\
 & fused two-stream        & $0.550\pm0.112$ & $0.528\pm0.081$\\
 & concat+MLP              & $0.549\pm0.089$ & $0.520\pm0.093$\\
\bottomrule
\end{tabular}
\end{table}

\begin{table}[h]\centering
\caption{Backbone (PLM) ablation on SAAINT-DB (10-fold, mean $\pm$ s.d.). ESM-2 is the strongest encoder
by all three metrics.}
\begin{tabular}{lccc}
\toprule
Backbone & Pearson $r$ & Spearman $\rho$ & RMSE (\pkd{})\\
\midrule
ESM-2      & $\mathbf{0.838\pm0.028}$ & $\mathbf{0.822\pm0.033}$ & $\mathbf{0.727\pm0.045}$\\
ProtBERT   & $0.819\pm0.025$ & $0.809\pm0.027$ & $0.770\pm0.047$\\
AntiBERTy  & $0.532\pm0.111$ & $0.517\pm0.117$ & $1.100\pm0.092$\\
ProGen2    & $0.412\pm0.036$ & $0.393\pm0.030$ & $1.165\pm0.053$\\
\bottomrule
\end{tabular}
\end{table}

\begin{table}[h]\centering
\caption{External benchmarks: AbAffinity vs.\ MVSF-AB (Pearson $r$ and RMSE in \pkd{}).}
\begin{tabular}{lcccc}
\toprule
& \multicolumn{2}{c}{Pearson $r$} & \multicolumn{2}{c}{RMSE (\pkd{})}\\
\cmidrule(lr){2-3}\cmidrule(lr){4-5}
Benchmark & MVSF-AB & AbAffinity & MVSF-AB & AbAffinity\\
\midrule
SAbDab       & 0.491 & $\mathbf{0.593}$ & 1.839 & $\mathbf{1.326}$\\
AB-Bind      & 0.739 & $\mathbf{0.781}$ & 1.905 & $\mathbf{1.288}$\\
SKEMPI 2.0   & 0.671 & $\mathbf{0.722}$ & 1.513 & $\mathbf{1.021}$\\
held-out     & 0.467 & $\mathbf{0.551}$ & 1.447 & $\mathbf{1.339}$\\
\bottomrule
\end{tabular}
\end{table}

\begin{table}[h]\centering
\caption{Few-shot fine-tuning (within-target Spearman) at $0\%$ (zero-shot) and $30\%$ labelled variants.}
\begin{tabular}{lcc}
\toprule
Target & zero-shot & fine-tuned (30\%)\\
\midrule
1mlc         & 0.17 & $\mathbf{0.42}$\\
trastuzumab  & 0.14 & $\mathbf{0.30}$\\
4fqi         & 0.47 & $\mathbf{0.95}$\\
\bottomrule
\end{tabular}
\end{table}

\section{Mutational \ddg{} supplementary tables}

\begin{table}[h]\centering
\caption{S1131 \ddg{} generalization: AbAffinity vs.\ frozen-ESM-2 regressors and a parameter-matched
LSTM on \emph{identical folds} ($z$-pooled Pearson, mean $\pm$ s.d.\ over 3 seeds). The affinity anchor's
advantage grows as the split becomes stricter.}
\begin{tabular}{lccc}
\toprule
Model & random & complex-disjoint & antigen-disjoint\\
\midrule
AbAffinity-\ddg{} (ours) & $0.801\pm0.004$ & $\mathbf{0.713\pm0.017}$ & $\mathbf{0.675\pm0.019}$\\
affinity anchor only     & $0.722$ & $0.687$ & $0.650$\\
XGBoost (ESM-2)          & $0.851$ & $0.580$ & $0.458$\\
RandomForest (ESM-2)     & $0.842$ & $0.540$ & $0.414$\\
MLP (ESM-2)              & $0.819$ & $0.363$ & $0.367$\\
LSTM (param-matched)     & $0.60$  & $0.161$ & $0.061$\\
\bottomrule
\end{tabular}
\end{table}

\begin{table}[h]\centering
\caption{S1131 random 10-fold CV: AbAffinity-\ddg{} vs.\ published predictors on the same split.
AbAffinity leads all sequence-only methods; fine-tuned ProtAttBA and structure-based DDGPred are higher.}
\begin{tabular}{lc}
\toprule
Method & Pearson $r$\\
\midrule
AbAffinity-\ddg{} (frozen sequence, ours) & $\mathbf{0.80}$\\
AttABseq (sequence)      & 0.66\\
DeepEP-PPI (sequence)    & 0.41\\
TransPPI (sequence)      & 0.38\\
PIPR (sequence)          & 0.33\\
BeAtMuSiC (sequence)     & 0.29\\
\midrule
ProtAttBA (fine-tuned sequence) & 0.84\\
DDGPred (structure)             & 0.92\\
\bottomrule
\end{tabular}
\end{table}

\begin{table}[h]\centering
\caption{Anchor-formulation ablation (S1131, $z$-pooled Pearson, mean $\pm$ s.d.). The three
thermodynamically-equivalent parameterizations of the affinity anchor are statistically identical; the
\pkd{}/\ddg{} forms carry lower seed variance. We adopt the \pkd{} form (main model).}
\begin{tabular}{lccc}
\toprule
Anchor & random & complex-disjoint & antigen-disjoint\\
\midrule
\pkd{} difference $s(\pkd^{wt}-\pkd^{mut})$ (main) & $0.801\pm0.004$ & $0.713\pm0.017$ & $0.675\pm0.019$\\
cosine difference $s(\cos^{wt}-\cos^{mut})$        & $0.803\pm0.008$ & $0.700\pm0.010$ & $0.670\pm0.039$\\
$\Delta G$ difference $s(\Delta G^{mut}-\Delta G^{wt})$ & $0.792\pm0.006$ & $0.726\pm0.008$ & $0.682\pm0.018$\\
\bottomrule
\end{tabular}
\end{table}

\begin{table}[h]\centering
\caption{Large-effect detection on S1131 (mutations with $|\ddg|>2$~kcal/mol): direction accuracy and
AUROC for identifying affinity-enhancing substitutions ($\ddg<0$).}
\begin{tabular}{lccc}
\toprule
Metric & random & complex-disjoint & antigen-disjoint\\
\midrule
direction accuracy       & 0.97 & 0.96 & 0.94\\
AUROC (enhancing vs not) & 0.98 & 0.95 & 0.95\\
\bottomrule
\end{tabular}
\end{table}

\section{Datasets and splits}
\textbf{SAAINT-DB} (affinity): 2{,}575 non-redundant complexes / 1{,}271 PDBs (1{,}913 paired antibodies,
662 nanobodies). \textbf{Splits}: random 10-fold; antigen-cold 10-fold (every PDB/antigen group assigned to
one fold). \textbf{External}: MVSF-AB release --- SAbDab, AB-Bind, SKEMPI~2.0, held-out. \textbf{\ddg{}}:
S1131 (1{,}130 single-point interface mutations, 112 complexes), AB-Bind AB645/AB1101, and general
protein--protein SKEMPI~2.0 for transfer. \ddg{} split ladder: random (record-level), complex-disjoint
(\texttt{GroupKFold} by PDB / \texttt{fold\_complex5}), antigen-disjoint (\texttt{GroupKFold} by antigen /
\texttt{fold\_antigen5}). Baselines are run on the \emph{identical} fold columns as AbAffinity.

\section{Interpretability details}
Integrated Gradients uses a zero-vector baseline and 50 interpolation steps; per-residue scores are summed
over the 1{,}280 ESM-2 dimensions and normalized within each chain. Ground-truth interfaces are taken from
the deposited crystal structures. On 1VFB (D1.3--lysozyme), AbAffinity recovers 19/22 antibody paratope
residues (heavy 10/12, light 9/10) at attribution threshold $0.4$, vs.\ 7/22 for a two-stream ablation;
aggregated over the five annotated complexes it recovers 56/73 paratope residues vs.\ 22/73 for a fused
baseline.

\end{document}
